% ============================================
% Chapter 6: Capture Environment and Nmap Simulations
% ============================================
% Condensed from Moulay Anas's main.tex (30+ pages to ~10 pages)
% Original: French, full detail; this chapter: English, essential content only

\chapter{Capture Environment and Nmap Simulations}

\section{Introduction}

This chapter describes the infrastructure and experimental setup used to generate the port scan traffic analyzed throughout the project. A controlled virtual laboratory was built to produce realistic reconnaissance scenarios, providing reproducible attack traces for both model training (via the CICIDS2017 dataset) and real-time system validation.

\section{Virtualization Environment}

The laboratory relies on Oracle VirtualBox. Two virtual machines are configured with distinct roles. Figure~\ref{fig:vbox-global} shows the VirtualBox manager with both laboratory machines.

\begin{figure}[H]
    \centering
    \includegraphics[width=0.95\linewidth]{image (1).png}
    \caption{VirtualBox manager overview showing both laboratory virtual machines.}
    \label{fig:vbox-global}
\end{figure}

\subsection{Attacker Machine: Kali Linux}

Kali Linux is configured with two network interfaces: a NAT interface for external access and an internal interface on the \texttt{lab-cyber} network for communicating directly with the target.

\begin{figure}[H]
    \centering
    \includegraphics[width=0.68\linewidth]{image (2).png}
    \caption{VirtualBox configuration of the Kali Linux attacker machine.}
    \label{fig:kali-vbox}
\end{figure}

\begin{table}[H]
\centering
\caption{Kali Linux virtual machine configuration}
\label{tab:kali-config}
\begin{tabular}{ll}
\toprule
\textbf{Property} & \textbf{Value} \\
\midrule
Machine name & \texttt{kali-linux-2026.1-virtualbox-amd64} \\
Declared OS & Debian 64-bit \\
Role & Attacker, source of Nmap scans \\
RAM & 2048 MB \\
CPUs & 2 \\
Storage & $\sim$80 GB (VDI) \\
Interface 1 & NAT (external access if needed) \\
Interface 2 & Internal network \texttt{lab-cyber} (scan traffic) \\
\bottomrule
\end{tabular}
\end{table}

\subsection{Target Machine: Ubuntu Victime}

Ubuntu Victime is placed on the same internal network as Kali to enable reconnaissance testing. It also retains a separate NAT interface, which is not used for scans.

\begin{figure}[H]
    \centering
    \includegraphics[width=0.68\linewidth]{image (3).png}
    \caption{VirtualBox configuration of the Ubuntu Victime target machine.}
    \label{fig:ubuntu-vbox}
\end{figure}

\begin{table}[H]
\centering
\caption{Ubuntu Victime virtual machine configuration}
\label{tab:ubuntu-config}
\begin{tabular}{ll}
\toprule
\textbf{Property} & \textbf{Value} \\
\midrule
Machine name & Ubuntu Victime \\
Declared OS & Ubuntu 25.04, 64-bit \\
Role & Target, receives Nmap scans \\
RAM & 4096 MB \\
CPUs & 2 \\
Storage & $\sim$25 GB (VDI) \\
Interface 1 & NAT (VM internet access) \\
Interface 2 & Internal network \texttt{lab-cyber} (target interface) \\
\bottomrule
\end{tabular}
\end{table}

\subsection{Network Separation}

The configuration uses two interface types. The NAT interface allows machines to access the external network when needed. The \texttt{lab-cyber} internal network is used exclusively for communication between attacker and victim. This separation is important because it prevents test traffic from mixing with external traffic. In all simulations, Nmap reconnaissance packets traverse the internal network, not the NAT interface, keeping attack traffic confined to the virtual laboratory.

\section{Network Architecture}

\subsection{Topology}

The topology consists of two machines connected via VirtualBox's internal network \texttt{lab-cyber}, using the \texttt{192.168.100.0/24} address space.

\begin{figure}[H]
\centering
\begin{tikzpicture}[
    node distance=2.2cm,
    vm/.style={draw, rounded corners, thick, minimum width=4.2cm, minimum height=1.6cm, align=center, fill=LightGray},
    net/.style={draw, rounded corners, thick, minimum width=4.4cm, minimum height=1.3cm, align=center, fill=LightBlue},
    arrow/.style={-{Latex[length=3mm]}, thick}
]
    \node[vm] (kali) {\textbf{Kali Linux}\\eth1: 192.168.100.10/24\\Role: attacker};
    \node[net, right=2.2cm of kali] (lab) {\textbf{lab-cyber}\\192.168.100.0/24};
    \node[vm, right=2.2cm of lab] (ubuntu) {\textbf{Ubuntu Victime}\\enp0s8: 192.168.100.20/24\\Role: target};
    \draw[thick] (kali) -- (lab);
    \draw[thick] (lab) -- (ubuntu);
\end{tikzpicture}
\caption{Logical topology of the test laboratory.}
\label{fig:topologie-logique}
\end{figure}

\subsection{Addressing Plan}

\begin{table}[H]
\centering
\caption{Addressing plan and interface roles}
\label{tab:adressage}
\begin{tabular}{llll}
\toprule
\textbf{Machine} & \textbf{Interface} & \textbf{IP Address} & \textbf{Role} \\
\midrule
Kali Linux & eth1 & 192.168.100.10/24 & Nmap scan source \\
Ubuntu Victime & enp0s8 & 192.168.100.20/24 & Nmap scan target \\
Ubuntu Victime & enp0s3 & 10.0.2.15/24 & VM internet (NAT), unused for scans \\
\bottomrule
\end{tabular}
\end{table}

\noindent The NAT interfaces remain separate from the attack path. Scan traffic flows exclusively through the internal \texttt{lab-cyber} network, keeping test traffic isolated from external networks.

\section{Nmap Scan Scenarios}

Ten Nmap scan scenarios were executed from Kali Linux toward Ubuntu Victime, covering a range of reconnaissance behaviors from stealthy to aggressive. Each scan produced three output files (\texttt{.nmap}, \texttt{.xml}, \texttt{.gnmap}) via the \texttt{-oA} option, yielding 30 files total.

\subsection{Scan Commands}

\begin{lstlisting}[caption={Nmap commands executed from Kali Linux}]
nmap -sS -p 1-1000 -oA ~/scans/01_syn_scan 192.168.100.20
nmap -sS -p- -oA ~/scans/02_full_tcp_scan 192.168.100.20
nmap -sV -p 1-1000 -oA ~/scans/03_service_detection 192.168.100.20
nmap -O -oA ~/scans/04_os_detection 192.168.100.20
nmap -A -oA ~/scans/05_aggressive_scan 192.168.100.20
nmap -sU --top-ports 50 -oA ~/scans/06_udp_scan 192.168.100.20
nmap -sX -p 1-1000 -oA ~/scans/07_xmas_scan 192.168.100.20
nmap -sS -p 1-1000 --scan-delay 500ms --max-rate 5 \
    -oA ~/scans/08_slow_scan 192.168.100.20
nmap -sS -T4 -p 1-1000 -oA ~/scans/09_fast_scan 192.168.100.20
nmap -sS -T5 -p 1-1000 -oA ~/scans/10_very_fast_scan 192.168.100.20
\end{lstlisting}

\subsection{Summary of Results}

\begin{longtable}{p{3.2cm}p{2.2cm}p{7.5cm}}
\caption{Summary of the ten Nmap scans} \label{tab:synthese-scans} \\
\toprule
\textbf{Scan} & \textbf{Duration} & \textbf{Key Result} \\
\midrule
\endfirsthead
\toprule
\textbf{Scan} & \textbf{Duration} & \textbf{Key Result} \\
\midrule
\endhead
TCP SYN scan & 4.67 s & 22/tcp open SSH; 80/tcp open HTTP \\
Full TCP scan & 6.24 s & 22/tcp open SSH; 80/tcp open HTTP \\
Service detection & 10.91 s & OpenSSH 10.2p1; Apache httpd 2.4.66 \\
OS detection & 15.96 s & No exact OS match; network distance: 1 hop \\
Aggressive scan & 22.41 s & Services + HTTP title + traceroute \\
UDP scan & 20.47 s & No confirmed open UDP ports; 30 \texttt{open|filtered} \\
XMAS scan & 5.86 s & 22/tcp and 80/tcp appear \texttt{open|filtered} \\
Slow scan (low-and-slow) & 505.83 s & 22/tcp open SSH; 80/tcp open HTTP \\
Fast scan (T4) & 4.66 s & 22/tcp open SSH; 80/tcp open HTTP \\
Very fast scan (T5) & 4.66 s & 22/tcp open SSH; 80/tcp open HTTP \\
\bottomrule
\end{longtable}

\subsection{Open Services on the Target}

TCP scans consistently identified two open services on Ubuntu Victime:

\begin{table}[H]
\centering
\caption{Open services detected on Ubuntu Victime}
\label{tab:services}
\begin{tabular}{lll}
\toprule
\textbf{Port} & \textbf{Service} & \textbf{Details} \\
\midrule
22/tcp & ssh & OpenSSH 10.2p1 Ubuntu 2ubuntu3.2 \\
80/tcp & http & Apache httpd 2.4.66; default Apache page \\
\bottomrule
\end{tabular}
\end{table}

\section{Relevance to the Detection System}

The scan scenarios serve multiple purposes in the AEGIS Entropy project:

\begin{enumerate}
    \item \textbf{Attack diversity}: The ten scans cover SYN, full TCP, service detection, OS detection, aggressive, UDP, XMAS, slow, fast, and very fast profiles --- providing a broad spectrum of reconnaissance behaviors.
    \item \textbf{Low-and-slow simulation}: The slow scan (505.83 s with \texttt{--scan-delay 500ms --max-rate 5}) is particularly relevant for testing detection of stealthy reconnaissance.
    \item \textbf{AI validation}: The XML output files are structured for automated parsing and can feed the ML pipeline for real-time classification.
    \item \textbf{Reactivity measurement}: The gap between scan start time and alert generation defines the system's detection latency:
    \[
    T_{\text{reactivity}} = T_{\text{alert}} - T_{\text{scan start}}
    \]
\end{enumerate}

\section{Conclusion}

The virtual laboratory provides a controlled, reproducible environment for generating and analyzing port scan traffic. Ten Nmap scan scenarios --- spanning stealthy to aggressive --- were successfully executed against the Ubuntu Victime target. The 30 output files (3 formats $\times$ 10 scans) constitute the capture layer's dataset, which integrates with the ML detection pipeline described in earlier chapters and the real-time monitoring dashboard presented in the following chapter.
